\section{EXPERIMENTAL METHOD}
\subsection{Content}
Conducting the experimental drying at three different temperatures of calorifer: $50^\circ C$, $60^\circ C$, and $70^\circ C$.

\subsection{Experiments}
\subsubsection{Observation of dryer}
Observe the drying apparatus and compare with the diagram before practice.

\subsubsection{Preparation}
\begin{itemize}
    \item Weighing the initial dry filter paper ($G_0$) using the mass balance of the dryer.
    \item Sufficiently humidifying the filter paper, then determining the weight ($G$).
    \item Preparing and setting zero for the timer.
\end{itemize}

\subsubsection{Dryer checking}
\begin{itemize}
    \item Close the door of the drying oven.
    \item Check the valve of air inlet and outlet doors.
    \item Add water into the water gourd of the wet-bulb thermometer.
\end{itemize}

\subsubsection{Start of the dryer}
\begin{itemize}
    \item \textbf{Blower}: Switch on for loading humid air and feeding to the calorifer.
    \item \textbf{Calorifer}: Switch on. Turn on the first resistor for $50^\circ C$, and the second for higher temperatures.
\end{itemize}

\subsection{Collecting data in an experiment}
\subsubsection{The data need collecting}
The weight of filter paper before, after and during drying process, dry and wet-bulb temperature, drying time.
\subsubsection{Reporting the weight by mass balance}
\begin{itemize}
    \item The weight ($g$): reporting the small value on the mass balance
    \item Temperature ($^\circ C$): reporting the value on the temperature indicator
\end{itemize}

\subsection{Accomplishing}
\begin{enumerate}
    \item Turn off all dry resistors of the calorifer.
    \item Switch off the blower after shutting down the calorifer.
    \item Remove filter paper from the drying oven.
    \item Switch off the circuit breaker (CP) of the dryer.
\end{enumerate}
